\chapter{Przestrzeń probabilistyczna i miara prawdopodobieństwa}
\label{ch:probability}

\section{Zdarzenia i \texorpdfstring{$\sigma$}{sigma}-ciała}
W doświadczeniu losowym najpierw występuje pewien wynik
$\omega\in\Omega$. Z reguły nie chcemy jednak przyznawać
prawdopodobieństwa pojedynczym wynikom, lecz formułować zdania,
na przykład „suma oczek jest parzysta” lub „ciąg wyników ma
nieskończenie wiele jedynek”. Takie zdania odpowiadają podzbiorom
przestrzeni wyników. Operacje logiczne „nie”, „lub”, „i” odpowiadają
dopełnieniu, sumie i przecięciu zbiorów.

Dlaczego nie zawsze bierzemy wszystkie podzbiory $\Omega$?
Dla przestrzeni nieprzeliczalnych próba przypisania długości
każdemu podzbiorowi odcinka z zachowaniem standardowych
własności miary prowadzi do niemożliwości. Dlatego z góry
ustalamy rodzinę $\mathcal F$ dopuszczalnych zdarzeń.
Zamknięcie na operacje \emph{przeliczalne} pozwoli mówić
między innymi o zdarzeniu „co najmniej jedno z $A_n$ zachodzi”
i o zdarzeniach granicznych.

\begin{definicja}
Niech $\Omega$ będzie niepustym zbiorem. Rodzinę $\mathcal F\subseteq2^\Omega$
nazywamy $\sigma$-ciałem, gdy (i) $\Omega\in\mathcal F$,
(ii) z $A\in\mathcal F$ wynika $\Omega\setminus A\in\mathcal F$,
(iii) dla każdego ciągu $(A_n)_{n\geq1}\subseteq\mathcal F$ zachodzi
$\bigcup_{n\geq1}A_n\in\mathcal F$.
Parę $(\Omega,\mathcal F)$ nazywamy przestrzenią mierzalną,
a elementy $\mathcal F$ zdarzeniami.
\end{definicja}

\begin{stwierdzenie}
Jeżeli $\mathcal F$ jest $\sigma$-ciałem na $\Omega$, to
$\varnothing\in\mathcal F$ oraz $\bigcap_{n\geq1}A_n\in\mathcal F$
dla każdego ciągu $A_n\in\mathcal F$. Rodzina $\mathcal F$ jest
również zamknięta na skończone sumy, przekroje i różnice zbiorów.
\end{stwierdzenie}
\begin{proof}
Ponieważ $\Omega\in\mathcal F$ i $\mathcal F$ jest zamknięte
na dopełnienia, to $\varnothing=\Omega^c\in\mathcal F$.
Jeśli $A_1,A_2,\ldots\in\mathcal F$, to każde
$A_n^c$ również należy do $\mathcal F$. Wobec tego
$\bigcup_{n\geq1}A_n^c\in\mathcal F$, a przez ponowne
wzięcie dopełnienia i prawo de Morgana dostajemy
\[
\bigcap_{n\geq1}A_n
=\left(\bigcup_{n\geq1}A_n^c\right)^c\in\mathcal F.
\]
Sumy i przekroje skończone są przypadkami przeliczalnymi:
do skończonej rodziny dopisujemy na końcu zbiory puste
lub całe $\Omega$, odpowiednio. Różnicę zapisujemy jako
$A\setminus B=A\cap B^c$, co dowodzi pozostałych własności.
\end{proof}

\begin{definicja}
Dla przestrzeni topologicznej $E$ przez $\mathcal B(E)$ oznaczamy
najmniejsze $\sigma$-ciało zawierające wszystkie zbiory otwarte w $E$;
nazywamy je $\sigma$-ciałem borelowskim.
\end{definicja}

\begin{przyklad}
Dla zbioru skończonego $\Omega$ przykładem jest $2^\Omega$.
Dla dowolnego $\Omega$ przykładem jest $\{\varnothing,\Omega\}$.
Jeżeli $A$ jest niepustym właściwym podzbiorem $\Omega$, to
$\{\varnothing,A,\Omega\setminus A,\Omega\}$ jest $\sigma$-ciałem:
pozwala rozstrzygać jedynie, w której z dwóch grup leży wynik.
Na $\R$ najczęściej używamy $\mathcal B(\R)$, a nie całego $2^\R$.
\end{przyklad}

\section{Aksjomaty Kołmogorowa i ich konsekwencje}
Dla skończonego zbioru wyników wystarczy podać wagi punktów,
które sumują się do jedności. W przestrzeniach nieskończonych
potrzebujemy zasady zgodności również dla przeliczalnych sum
rozłącznych zdarzeń. Stąd przeliczalna addytywność: jeśli zdarzenie
rozłożono na przeliczalnie wiele wykluczających się przypadków,
jego masa musi być sumą ich mas. Nie zakładamy wszystkich znanych
wzorów rachunku prawdopodobieństwa, lecz \emph{wyprowadzamy} je
z dwóch warunków: normalizacji i przeliczalnej addytywności.

\begin{definicja}
Niech $(\Omega,\mathcal F)$ będzie przestrzenią mierzalną.
Miarą probabilistyczną na $\mathcal F$ nazywamy funkcję
$\Prob:\mathcal F\to[0,1]$, dla której $\Prob(\Omega)=1$
oraz dla każdego ciągu parami rozłącznych zdarzeń $A_n\in\mathcal F$
zachodzi
\begin{equation}\label{eq:sigma_add}
\Prob\!\left(\bigcup_{n=1}^\infty A_n\right)
=\sum_{n=1}^\infty\Prob(A_n).
\end{equation}
Trójkę $(\Omega,\mathcal F,\Prob)$ nazywamy przestrzenią probabilistyczną.
\end{definicja}

\begin{twierdzenie}[Podstawowe własności miary probabilistycznej]
Niech $(\Omega,\mathcal F,\Prob)$ będzie przestrzenią probabilistyczną.
Dla wszystkich $A,B\in\mathcal F$ oraz dowolnego ciągu
$(A_n)_{n\geq1}\subseteq\mathcal F$ zachodzą:
\begin{enumerate}[label=(\roman*)]
\item $\Prob(\varnothing)=0$ i $\Prob(A^c)=1-\Prob(A)$;
\item jeśli $A\subseteq B$, to $\Prob(A)\leq\Prob(B)$ oraz
      $\Prob(B\setminus A)=\Prob(B)-\Prob(A)$;
\item $\Prob(A\cup B)=\Prob(A)+\Prob(B)-\Prob(A\cap B)$;
\item $\Prob(\bigcup_{n\geq1}A_n)\leq\sum_{n\geq1}\Prob(A_n)$.
\end{enumerate}
\end{twierdzenie}
\begin{proof}
\emph{(i)} Z przeliczalnej addytywności zastosowanej do
rozłącznego ciągu $\Omega,\varnothing,\varnothing,\ldots$
otrzymujemy
$1=\Prob(\Omega)=1+\sum_{n\geq2}\Prob(\varnothing)$.
Nieujemność każdego składnika wymusza
$\Prob(\varnothing)=0$. Następnie
$\Omega=A\mathbin{\dot\cup}A^c$ daje
$1=\Prob(A)+\Prob(A^c)$.

\emph{(ii)} Jeśli $A\subseteq B$, to
$B=A\mathbin{\dot\cup}(B\setminus A)$.
Z addytywności mamy
$\Prob(B)=\Prob(A)+\Prob(B\setminus A)$, a ostatni
składnik jest nieujemny. Jednocześnie otrzymujemy wzór
na masę różnicy.

\emph{(iii)} Zapiszmy
$A\cup B=A\mathbin{\dot\cup}(B\setminus A)$ oraz
$B=(B\setminus A)\mathbin{\dot\cup}(A\cap B)$.
Stosując skończoną addytywność do obu równości i
eliminując $\Prob(B\setminus A)$, otrzymujemy
$\Prob(A\cup B)=\Prob(A)+\Prob(B)-\Prob(A\cap B)$.

\emph{(iv)} Dla ogólnych $A_1,A_2,\ldots$ nie możemy
od razu użyć addytywności, bo zbiory mogą się przecinać.
Budujemy więc ich rozłącznienie: $B_1=A_1$ oraz
$B_n=A_n\setminus\bigcup_{j=1}^{n-1}A_j$ dla $n\geq2$.
Wtedy zbiory $B_n$ są parami rozłączne,
$\bigcup_n B_n=\bigcup_n A_n$ i $B_n\subseteq A_n$.
Zatem dzięki monotoniczności i przeliczalnej addytywności
\[
\Prob\left(\bigcup_n A_n\right)
=\sum_n\Prob(B_n)\leq\sum_n\Prob(A_n).
\]
\end{proof}

\begin{twierdzenie}[Ciągłość miary]
Niech $(\Omega,\mathcal F,\Prob)$ będzie przestrzenią probabilistyczną.
Jeśli $A_n\uparrow A$, to $\Prob(A_n)\uparrow\Prob(A)$.
Jeśli $A_n\downarrow A$, to $\Prob(A_n)\downarrow\Prob(A)$.
Tutaj $A_n\uparrow A$ oznacza $A_n\subseteq A_{n+1}$ i
$A=\bigcup_{n\geq1}A_n$, a $A_n\downarrow A$ oznacza
$A_{n+1}\subseteq A_n$ i $A=\bigcap_{n\geq1}A_n$.
\end{twierdzenie}
\begin{proof}
Załóżmy najpierw, że $A_n\uparrow A$.
Wprowadźmy rozłączne zdarzenia
$D_1=A_1$ i $D_n=A_n\setminus A_{n-1}$ dla $n\geq2$.
Ponieważ ciąg jest rosnący, $A_k$ jest dokładnie sumą
$D_1,\ldots,D_k$, a $A$ jest sumą wszystkich $D_n$.
Z addytywności
$\Prob(A_k)=\sum_{n=1}^k\Prob(D_n)$ oraz
$\Prob(A)=\sum_{n=1}^\infty\Prob(D_n)$.
Granice sum częściowych nieujemnego szeregu dają
$\lim_{k\to\infty}\Prob(A_k)=\Prob(A)$.

Jeżeli $A_n\downarrow A$, to ciąg $A_n^c$ jest rosnący,
a z praw de Morgana
$\bigcup_n A_n^c=(\bigcap_n A_n)^c=A^c$.
Pierwsza część dowodu daje
$\Prob(A_n^c)\uparrow\Prob(A^c)$. Korzystając ze wzoru
$\Prob(E^c)=1-\Prob(E)$ dla $E\in\mathcal F$,
otrzymujemy
$\Prob(A_n)=1-\Prob(A_n^c)\downarrow1-\Prob(A^c)
=\Prob(A)$.
\end{proof}

\begin{uwaga}
Warunek $\Prob(\Omega)=1$ jest istotny dla ciągłości z góry
bez dodatkowych ograniczeń: dla miary ogólnej potrzeba
$\mu(A_1)<\infty$. W przestrzeni probabilistycznej jest on
spełniony automatycznie.
\end{uwaga}

\begin{stwierdzenie}[Włączenia i wyłączenia dla prawdopodobieństw]
Dla każdej przestrzeni probabilistycznej i zdarzeń $A_1,\ldots,A_m$
zachodzi
\[
\Prob\!\left(\bigcup_{i=1}^m A_i\right)
=\sum_{\varnothing\ne I\subseteq\{1,\ldots,m\}}
(-1)^{|I|+1}\Prob\!\left(\bigcap_{i\in I}A_i\right).
\]
\end{stwierdzenie}
\begin{proof}
Dowodzimy twierdzenia indukcyjnie względem liczby zbiorów.
Dla $m=1$ jest oczywiste, a dla $m=2$ jest to
wyprowadzona wcześniej tożsamość
$\Prob(A\cup B)=\Prob(A)+\Prob(B)-\Prob(A\cap B)$.
Załóżmy, że wzór obowiązuje dla $m-1$ zbiorów.
Oznaczmy $U=\bigcup_{i=1}^{m-1}A_i$. Wtedy
\[
\Prob(U\cup A_m)
=\Prob(U)+\Prob(A_m)-\Prob(U\cap A_m).
\]
Do pierwszego składnika stosujemy założenie indukcyjne
dla $A_1,\ldots,A_{m-1}$. Do ostatniego stosujemy
je dla zbiorów $A_1\cap A_m,\ldots,A_{m-1}\cap A_m$,
ponieważ
$U\cap A_m=\bigcup_{i=1}^{m-1}(A_i\cap A_m)$.
Po odjęciu tej drugiej sumy każdy przekrój indeksowany
zbiorem $I$ bez $m$ ma znak $(-1)^{|I|+1}$,
a każdy przekrój zawierający $m$ otrzymuje znak
$-(-1)^{|I\setminus\{m\}|+1}=(-1)^{|I|+1}$.
Pozostaje jeszcze $\Prob(A_m)$ odpowiadające $I=\{m\}$.
Są to wszystkie niepuste zbiory indeksów, co kończy indukcję.
\end{proof}


\begin{przyklad}[Prawie pewne zdarzenie nie musi zachodzić zawsze]
Niech $U$ ma rozkład jednostajny na $(0,1)$.
Wtedy zdarzenie $A=\{U\neq1/2\}$ ma prawdopodobieństwo
jeden, mimo że na przestrzeni wyników $\Omega=(0,1)$
nie jest równe całemu $\Omega$. Mówimy, że $A$ zachodzi
\emph{prawie na pewno}. Z drugiej strony zdarzenie
$\{U=1/2\}$ jest logicznie możliwe w tym modelu,
lecz ma prawdopodobieństwo zero.
\end{przyklad}


\section{Dodatkowe własności przestrzeni probabilistycznej}
Często przy formułowaniu twierdzeń używa się wyrażeń
„prawie na pewno”, „poza zbiorem zerowym” i „zdarzenie
niemal niemożliwe”. Pierwsze dwa określenia mają ścisłe
znaczenie matematyczne; trzecie jest tylko nieformalnym
opisem zdarzenia o małym, lecz dodatnim prawdopodobieństwie.

\begin{definicja}
Niech $(\Omega,\mathcal F,\Prob)$ będzie przestrzenią
probabilistyczną i $A\in\mathcal F$. Mówimy, że $A$
zachodzi \emph{prawie na pewno}, gdy $\Prob(A)=1$.
Mówimy, że $A$ jest zdarzeniem \emph{zerowym}, gdy
$\Prob(A)=0$. Własność określona dla punktów $\omega\in\Omega$
zachodzi prawie na pewno, jeśli zbiór punktów, na których
nie zachodzi, należy do $\mathcal F$ i ma prawdopodobieństwo zero.
\end{definicja}

\begin{stwierdzenie}[Zbiory zerowe i prawie pewne]
Niech $(\Omega,\mathcal F,\Prob)$ będzie przestrzenią
probabilistyczną. Jeśli $N_n\in\mathcal F$ i
$\Prob(N_n)=0$ dla każdego $n\geq1$, to
$\Prob(\bigcup_{n\geq1}N_n)=0$. Jeśli $A_n\in\mathcal F$
i $\Prob(A_n)=1$ dla każdego $n\geq1$, to
$\Prob(\bigcap_{n\geq1}A_n)=1$.
\end{stwierdzenie}
\begin{proof}
Z przeliczalnej podaddytywności otrzymujemy
$0\leq\Prob(\bigcup_nN_n)\leq\sum_n\Prob(N_n)=0$,
a więc pierwsza równość zachodzi. Dla drugiej zauważmy,
że $\Prob(A_n^c)=1-\Prob(A_n)=0$. Pierwsza część
dowodu daje
$\Prob(\bigcup_n A_n^c)=0$. Z praw de Morgana
$\bigcap_n A_n=(\bigcup_n A_n^c)^c$, dlatego
$\Prob(\bigcap_n A_n)=1$.
\end{proof}

\begin{przyklad}[Nieskończony ciąg rzutów i zdarzenie zerowe]
Niech $\Omega=\{0,1\}^{\mathbb N}$ z miarą niezależnych
uczciwych rzutów monetą. Dla $m\geq1$ niech $C_m$
oznacza zdarzenie, że pierwszych $m$ wyników to same orły.
Wtedy $C_{m+1}\subseteq C_m$ i $\Prob(C_m)=2^{-m}$.
Zdarzenie $C=\bigcap_{m\geq1}C_m$ oznacza nieskończony
ciąg samych orłów. Na mocy ciągłości miary z góry
$\Prob(C)=\lim_{m\to\infty}2^{-m}=0$.
Jednak $C$ nie jest puste: zawiera jeden konkretny wynik,
mianowicie nieskończony ciąg orłów. To ważny przykład,
w którym logicznie dopuszczalna realizacja ma masę zero.
\end{przyklad}

\begin{stwierdzenie}[Nierówność Boole'a dla zdarzeń]
Niech $A_1,\ldots,A_m$ będą zdarzeniami w przestrzeni
probabilistycznej, $m\geq1$. Wtedy
$\Prob(\bigcap_{i=1}^m A_i)
\geq 1-\sum_{i=1}^m\Prob(A_i^c)$.
\end{stwierdzenie}
\begin{proof}
Ponieważ $(\bigcap_i A_i)^c=\bigcup_i A_i^c$,
z dopełnienia i przeliczalnej podaddytywności otrzymujemy
$1-\Prob(\bigcap_i A_i)
=\Prob(\bigcup_i A_i^c)\leq\sum_i\Prob(A_i^c)$.
Przenosząc składniki, otrzymujemy żądaną nierówność.
Jest ona przydatna, gdy chcemy wykazać, że wiele
pożądanych własności zachodzi jednocześnie z dużym
prawdopodobieństwem, choć nie znamy zależności między nimi.
\end{proof}

\section{Niezależność zdarzeń}
Dwa zdarzenia mogą zachodzić jednocześnie, ale informacja
o jednym z nich nie musi zmieniać prawdopodobieństwa drugiego.
Tę sytuację formalizuje niezależność. Ponieważ nie chcemy
opierać definicji wyłącznie na warunkowaniu (które wymaga dodatniego
prawdopodobieństwa warunku), jako definicję podstawową przyjmujemy
równość dotyczącą przecięcia. Dla trzech i więcej zdarzeń
niezależność par nie wystarcza: zależność może być widoczna
dopiero po jednoczesnym uwzględnieniu wszystkich zdarzeń.

\begin{definicja}
Zdarzenia $A,B\in\mathcal F$ nazywamy niezależnymi, jeśli
$\Prob(A\cap B)=\Prob(A)\Prob(B)$.
Rodzinę zdarzeń $(A_i)_{i\in I}$ nazywamy wzajemnie niezależną,
jeśli dla każdej skończonej rodziny \emph{różnych} indeksów
$i_1,\ldots,i_m\in I$, $m\geq1$, zachodzi
$\Prob(\bigcap_{j=1}^m A_{i_j})=\prod_{j=1}^m\Prob(A_{i_j})$.
Rodzinę nazywamy parami niezależną, jeśli warunek ten obowiązuje
tylko dla par różnych indeksów.
\end{definicja}

\section{Zdarzenia zachodzące nieskończenie często}
W nieskończonym ciągu prób powstaje nowe pytanie: czy sukces
pojawi się choć raz, czy pojawi się po dowolnie późnym momencie,
czy może od pewnego miejsca będzie występował zawsze? Są to
trzy różne zdarzenia. Symbol $\limsup A_n$ oznacza, że
\emph{dla każdego} momentu $m$ znajdzie się późniejszy sukces.
Pierwszy lemat Borela--Cantellego wyraża zasadę, że zbyt mała
łączna suma szans sukcesów wyklucza nieskończenie wiele ich
wystąpień prawie na pewno. Drugi pokazuje, że przy niezależności
rozbieżność tej sumy działa w przeciwnym kierunku.

\begin{definicja}
Dla ciągu zdarzeń $(A_n)_{n\geq1}$ definiujemy
\begin{align*}
\limsup_{n\to\infty}A_n
&:=\bigcap_{m=1}^\infty\bigcup_{n=m}^\infty A_n,\\
\liminf_{n\to\infty}A_n
&:=\bigcup_{m=1}^\infty\bigcap_{n=m}^\infty A_n.
\end{align*}
Pierwsze zdarzenie oznacza, że zachodzi nieskończenie wiele zdarzeń $A_n$,
a drugie, że od pewnego miejsca zachodzą wszystkie zdarzenia $A_n$.
\end{definicja}


\begin{przyklad}[Kwantyfikatory ukryte w granicach zdarzeń]
Dla ustalonego $\omega\in\Omega$ zachodzi
\begin{align*}
\omega\in\limsup_{n\to\infty}A_n
&\iff (\forall m\geq1)(\exists n\geq m):\omega\in A_n,\\
\omega\in\liminf_{n\to\infty}A_n
&\iff (\exists m\geq1)(\forall n\geq m):\omega\in A_n.
\end{align*}
Jeżeli $A_n=\Omega$ dla parzystych $n$ i $A_n=\varnothing$
dla nieparzystych $n$, to $\limsup A_n=\Omega$, a
$\liminf A_n=\varnothing$. Jeśli zaś $A_n=A$ dla każdego $n$,
to obie granice są równe $A$. Zawsze mamy
$\liminf A_n\subseteq\limsup A_n$, ponieważ zdarzenie
zachodzące od pewnego momentu zawsze zachodzi nieskończenie często.
\end{przyklad}

\begin{twierdzenie}[Pierwszy lemat Borela--Cantellego]
Niech $(A_n)_{n\geq1}$ będzie ciągiem zdarzeń w przestrzeni
probabilistycznej $(\Omega,\mathcal F,\Prob)$.
Jeżeli $\sum_{n=1}^\infty\Prob(A_n)<\infty$, to
$\Prob(\limsup_{n\to\infty}A_n)=0$.
\end{twierdzenie}
\begin{proof}
Niech $L=\limsup_{n\to\infty}A_n$.
Jeżeli $\omega\in L$, to dla każdego ustalonego $m$
istnieje $n\geq m$ takie, że $\omega\in A_n$.
Zatem $L\subseteq U_m:=\bigcup_{n=m}^\infty A_n$.
Monotoniczność miary i nierówność sumy dają
\[
0\leq\Prob(L)\leq\Prob(U_m)
\leq\sum_{n=m}^\infty\Prob(A_n)
\qquad(m\geq1).
\]
Kluczowy jest ostatni wyraz: jest on ogonem 
\emph{zbieżnego szeregu nieujemnego}, toteż zmierza
do zera, gdy $m\to\infty$. Lewa strona nie zależy od $m$,
więc musi być zerem. Niezależność zdarzeń nie została
użyta w żadnym punkcie rozumowania.
\end{proof}


\begin{lemat}[Dopełnienia zdarzeń niezależnych]
Niech $A_1,\ldots,A_m$ będą wzajemnie niezależnymi
zdarzeniami w przestrzeni probabilistycznej, gdzie $m\geq1$.
Wówczas dla dowolnego $\varepsilon_1,\ldots,\varepsilon_m
\in\{0,1\}$ zdarzenia $B_i$, określone przez
$B_i=A_i$ gdy $\varepsilon_i=1$ i $B_i=A_i^c$
gdy $\varepsilon_i=0$, spełniają
$\Prob(\bigcap_{i=1}^m B_i)=\prod_{i=1}^m\Prob(B_i)$.
\end{lemat}
\begin{proof}
Dla dwóch zdarzeń korzystamy z rozłącznego rozkładu
$A=A\cap B\mathbin{\dot\cup}A\cap B^c$:
$\Prob(A\cap B^c)=\Prob(A)-\Prob(A\cap B)
=\Prob(A)(1-\Prob(B))$.
Dla dowolnej skończonej rodziny analogicznie zapisujemy
$C\cap A_j^c=C\setminus(C\cap A_j)$, gdzie $C$ jest
przecięciem już wybranych zdarzeń i ich dopełnień.
Zaczynając od wzoru na przecięcia samych $A_i$ i
kolejno zamieniając $A_j$ na $A_j^c$, uzyskujemy
czynnik $1-\Prob(A_j)$ w miejsce $\Prob(A_j)$.
Indukcja po liczbie dokonanych zamian daje tezę.
\end{proof}

\begin{twierdzenie}[Drugi lemat Borela--Cantellego]
Niech $(A_n)_{n\geq1}$ będzie ciągiem \emph{wzajemnie niezależnych}
zdarzeń w przestrzeni probabilistycznej $(\Omega,\mathcal F,\Prob)$.
Jeżeli $\sum_{n=1}^\infty\Prob(A_n)=\infty$, to
$\Prob(\limsup_{n\to\infty}A_n)=1$.
\end{twierdzenie}
\begin{proof}
Zauważmy najpierw, że
\[
(\limsup_{n\to\infty}A_n)^c
=\bigcup_{m=1}^\infty\bigcap_{n=m}^\infty A_n^c.
\]
Dopełnienie zdarzenia „$A_n$ występuje nieskończenie wiele
razy” oznacza bowiem „od pewnego momentu żadne $A_n$
już nie wystąpi”. Wystarczy więc pokazać, że dla każdego
$m$ zdarzenie $C_m:=\bigcap_{n=m}^\infty A_n^c$
ma prawdopodobieństwo zero.

Ustalmy $m$ i rozważmy skończone przekroje
$C_{m,N}=\bigcap_{n=m}^N A_n^c$ dla $N\geq m$.
Z niezależności i poprzedniego lematu otrzymujemy
\[
\Prob(C_{m,N})=\prod_{n=m}^N\bigl(1-\Prob(A_n)\bigr)
\leq \exp\!\left(-\sum_{n=m}^N\Prob(A_n)\right).
\]
Nierówność po prawej wynika ze
$1-u\leq e^{-u}$ dla $u\in[0,1]$, którą można otrzymać
na przykład z wypukłości funkcji wykładniczej.
Ponieważ cały szereg $\sum_n\Prob(A_n)$ jest rozbieżny,
jego ogon od ustalonego $m$ również ma sumy częściowe
dążące do $+\infty$. Prawa strona dąży więc do zera.
Zbiory $C_{m,N}$ maleją do $C_m$; dzięki ciągłości
miary z góry mamy
$\Prob(C_m)=\lim_{N\to\infty}\Prob(C_{m,N})=0$.
Na koniec przeliczalna podaddytywność daje
$\Prob((\limsup A_n)^c)
\leq\sum_{m=1}^\infty\Prob(C_m)=0$.
Oznacza to $\Prob(\limsup A_n)=1$.
\end{proof}

\begin{przyklad}[Niezależne próby]
W kolejnych niezależnych próbach zdarzenie sukcesu ma w próbie $n$
prawdopodobieństwo $p_n$. Jeśli $p_n=1/n^2$ dla $n\geq1$,
sukces następuje tylko skończenie wiele razy z prawdopodobieństwem jeden.
Jeśli $p_n=1/n$, sukces następuje nieskończenie wiele razy
z prawdopodobieństwem jeden. Pierwszy wniosek nie wymaga niezależności;
drugi jej wymaga.
\end{przyklad}

\begin{przyklad}[Dlaczego niezależność jest potrzebna]
Niech $A\in\mathcal F$ spełnia $\Prob(A)=1/2$ i niech $A_n=A$
dla wszystkich $n$. Wtedy $\sum_n\Prob(A_n)=\infty$,
lecz $\limsup A_n=A$, więc $\Prob(\limsup A_n)=1/2$.
Dla zdarzeń tych brak niezależności.
\end{przyklad}

\section{Konstrukcje przestrzeni probabilistycznych}
Abstrakcyjna definicja ma sens dopiero wtedy, gdy można zbudować
modele dla rzeczywistych doświadczeń. Zaczniemy od modelu dwóch
możliwości, potem przejdziemy do przestrzeni skończonych i
przeliczalnych, losowania punktu z odcinka oraz nieskończonego
ciągu rzutów monetą. Każda konstrukcja odpowiada na trzy pytania:
co może być wynikiem, jakie zdarzenia dopuszczamy i jak przypisujemy
im prawdopodobieństwa.


\begin{przyklad}[Wypadnie albo nie wypadnie: przestrzeń o prawdopodobieństwach $0$, $1/2$, $1$]
Załóżmy, że przy rzucie uczciwą monetą interesuje nas tylko,
czy wypadnie orzeł. Oznaczmy przez $o$ wynik „orzeł”, a przez
$r$ wynik „reszka”. Rozważamy
$\Omega=\{o,r\}$ oraz całe $\sigma$-ciało
$\mathcal F=2^\Omega=\{\varnothing,\{o\},\{r\},\Omega\}$.
Przy założeniu uczciwości monety definiujemy
$\Prob(\{o\})=\Prob(\{r\})=1/2$.
W konsekwencji pełna miara jest opisana przez tabelę
\[
\begin{array}{c|cccc}
A&\varnothing&\{o\}&\{r\}&\Omega\\ \hline
\Prob(A)&0&1/2&1/2&1
\end{array}
\]
Wartości prawdopodobieństw wszystkich dopuszczalnych zdarzeń
należą więc do zbioru $\{0,1/2,1\}$. Mamy cztery zdarzenia,
ale jedynie trzy różne wartości miary: dwóm różnym zdarzeniom
może odpowiadać to samo prawdopodobieństwo.

Sprawdźmy aksjomaty bez odwoływania się do intuicji. Suma
mas punktowych wynosi jeden, stąd $\Prob(\Omega)=1$.
Dla dowolnych parami rozłącznych $A_n\subseteq\Omega$ i
każdego $\omega\in\Omega$ punkt $\omega$ należy do
co najwyżej jednego $A_n$. Zatem
\[
\Prob\!\left(\bigcup_n A_n\right)
=\sum_{\omega\in\Omega}\frac12\ind{\{\omega\in\bigcup_n A_n\}}
=\sum_n\sum_{\omega\in A_n}\frac12
=\sum_n\Prob(A_n).
\]
Zamiana sum jest dopuszczalna, gdyż w tym przypadku tylko
skończenie wiele składników może być niezerowych.

\emph{Ważne:} sam opis „wypadnie albo nie wypadnie” nie oznacza,
że obie możliwości mają po $1/2$. Dla niesymetrycznego
mechanizmu z parametrem $p\in[0,1]$ możemy ustalić
$\Prob(\{o\})=p$ i $\Prob(\{r\})=1-p$.
Wtedy zbiór możliwych wartości miary to
$\{0,p,1-p,1\}$, który przy $p=1/2$ redukuje się
do $\{0,1/2,1\}$. Dwudzielność wyniku jest własnością
przestrzeni $\Omega$, a równość szans jest założeniem o miarze.
\end{przyklad}

\begin{stwierdzenie}[Charakteryzacja miar na przestrzeni dwupunktowej]
Niech $\Omega=\{o,r\}$, gdzie $o\ne r$, i niech
$\mathcal F=2^\Omega$. Każda miara probabilistyczna
$\Prob$ na $\mathcal F$ jest jednoznacznie określona przez
liczbę $p=\Prob(\{o\})\in[0,1]$ oraz ma postać
$\Prob(A)=p\ind{\{o\in A\}}+(1-p)\ind{\{r\in A\}}$
dla $A\in\mathcal F$. Odwrotnie, każda liczba $p\in[0,1]$
zadaje tym wzorem miarę probabilistyczną na $\mathcal F$.
\end{stwierdzenie}
\begin{proof}
Punkty tworzą rozłączny rozkład
$\Omega=\{o\}\mathbin{\dot\cup}\{r\}$, więc
$1=\Prob(\Omega)=\Prob(\{o\})+\Prob(\{r\})$.
Jeżeli pierwsza masa to $p$, druga musi wynosić $1-p$.
Każde zdarzenie jest sumą zawartych w nim singletonów, zatem
jego masa ma podaną postać. Na odwrót, dla ustalonego
$p\in[0,1]$ wzór ma wartości nieujemne, daje
$\Prob(\Omega)=p+(1-p)=1$ i jest przeliczalnie
addytywny, ponieważ każdy z dwóch punktów może leżeć
w co najwyżej jednym z parami rozłącznych zdarzeń.
\end{proof}

\begin{uwaga}[Prawdopodobieństwo zero a niemożliwość]
W modelu uczciwej monety jedynym zdarzeniem o masie zero jest
$\varnothing$, a więc prawdopodobieństwo zero oznacza tutaj
niemożliwość. Nie jest to ogólne prawo. W modelu jednostajnego
losowania $U$ z $[0,1]$ zdarzenie $\{U=1/2\}$ jest niepuste,
ale ma prawdopodobieństwo zero. Zatem implikacja
$A=\varnothing\Rightarrow\Prob(A)=0$ jest zawsze prawdziwa,
natomiast implikacja odwrotna może zawieść. Podobnie zdarzenie
o prawdopodobieństwie jeden nie musi być całym $\Omega$.
\end{uwaga}

\begin{przyklad}[Model skończony i model Laplace'a]
Niech $\Omega=\{\omega_1,\ldots,\omega_N\}$,
$\mathcal F=2^\Omega$ i niech $p_i\geq0$ oraz $\sum_{i=1}^N p_i=1$.
Wzór $\Prob(A)=\sum_{\omega_i\in A}p_i$ definiuje miarę
probabilistyczną: jest unormowana i addytywna przeliczalnie,
ponieważ w skończonym $\Omega$ tylko skończenie wiele rozłącznych
niepustych zbiorów może wystąpić. Przy $p_i=1/N$ otrzymujemy
$\Prob(A)=|A|/N$.
\end{przyklad}


\begin{przyklad}[Przestrzeń z czterema wynikami: dwa rzuty monetą]
Niech $\Omega=\{oo,or,ro,rr\}$ i $\mathcal F=2^\Omega$.
Zakładamy, że cztery wyniki mają masę $1/4$.
Wówczas zdarzenie $A$ „w pierwszym rzucie orzeł” to
$\{oo,or\}$, a zdarzenie $B$ „w drugim rzucie orzeł”
to $\{oo,ro\}$. Mamy $\Prob(A)=\Prob(B)=1/2$,
ale $\Prob(A\cap B)=1/4$, co odpowiada niezależności.
Zdarzenie $C=\{oo,rr\}$ również ma masę $1/2$,
choć oznacza inny warunek: dwa jednakowe wyniki.

Jeśli teraz obserwujemy jedynie, czy wyniki są jednakowe,
to możemy przejść do dwupunktowej przestrzeni wyników
$\Omega'=\{\text{jednakowe},\text{różne}\}$.
Zdarzenia $\{\text{jednakowe}\}$ i
$\{\text{różne}\}$ mają po $1/2$. Zmieniła się przestrzeń
wyników, ale prawdopodobieństwa zdarzeń wyrażalnych
przez tę obserwację pozostały takie same.
\end{przyklad}

\begin{przyklad}[Przestrzeń z niepełną informacją]
Niech $\Omega=\{1,2,3,4\}$ i
$\mathcal F=\{\varnothing,\{1,2\},\{3,4\},\Omega\}$.
Możemy pytać, czy wynik należy do pierwszej, czy drugiej
pary, lecz zdarzenie $\{1\}$ nie jest mierzalne.
Jeśli $\Prob(\{1,2\})=1/2$ oraz
$\Prob(\{3,4\})=1/2$, to prawdopodobieństwa wszystkich
\emph{dostępnych} zdarzeń ponownie należą do
$\{0,1/2,1\}$, mimo że sama przestrzeń wyników ma cztery
punkty. To pokazuje, że liczba wartości miary jest wspólną
cechą $\mathcal F$ oraz $\Prob$, nie samego $\Omega$.
\end{przyklad}

\begin{przyklad}[Model przeliczalny]
Niech $\Omega=\{1,2,\ldots\}$, $\mathcal F=2^\Omega$ i
$\Prob(A)=\sum_{n\in A}2^{-n}$. Ponieważ $\sum_{n\geq1}2^{-n}=1$,
jest to miara probabilistyczna. W szczególności
$\Prob(\{n\geq k\})=2^{-(k-1)}$ dla $k\geq1$.
\end{przyklad}

\begin{przyklad}[Jednostajny punkt odcinka]
Niech $\Omega=[0,1]$, $\mathcal F=\mathcal B([0,1])$,
a $\Prob$ będzie ograniczeniem jednowymiarowej miary Lebesgue'a
$\lambda_1$ do $\mathcal F$. Wtedy $\Prob([a,b])=b-a$
dla $0\leq a\leq b\leq1$. Wynikami elementarnymi są wszystkie
punkty odcinka, mimo że prawdopodobieństwo każdego pojedynczego
punktu wynosi zero. Dlatego skończonego wzoru Laplace'a nie można
tu stosować.
\end{przyklad}

\begin{przyklad}[Rzut monetą nieskończenie wiele razy]
Dla ustalonego $p\in[0,1]$ rozważamy $\Omega=\{0,1\}^{\N}$.
Zdarzenia cylindryczne określają skończoną liczbę współrzędnych.
Wymagamy, aby dla różnych indeksów $i_1,\ldots,i_m$ i
$a_j\in\{0,1\}$ zachodziło
\[
\Prob\{\omega:\omega_{i_j}=a_j,\ 1\leq j\leq m\}
=\prod_{j=1}^m p^{a_j}(1-p)^{1-a_j}.
\]
Na $\sigma$-ciele generowanym przez cylindryczne zdarzenia
istnieje dokładnie jedna miara probabilistyczna o tych wartościach;
jest to przeliczalny produkt miar Bernoulliego, uzyskiwany z twierdzenia
o istnieniu miary produktowej. Jego własności skończeniewymiarowe
opisują niezależne rzuty monetą, lecz pojedynczym wynikiem jest cały
nieskończony ciąg.
\end{przyklad}

\begin{uwaga}
Aksjomaty teorii prawdopodobieństwa nie przypisują wag dowolnym
podzbiorom nieprzeliczalnej przestrzeni $\Omega$. Zdarzenia muszą
należeć do wybranego $\sigma$-ciała. Ma to znaczenie choćby dla
nieborelowskich lub niemierzalnych w sensie Lebesgue'a podzbiorów $\R$.
\end{uwaga}

